\chapter*{Anexo B — Código fuente de la herramienta Python}
\addcontentsline{toc}{chapter}{Anexo B — Código fuente de la herramienta Python}
\label{anexo:python}

Este anexo recoge los fragmentos más representativos del código fuente de la
herramienta de diagnóstico implementada en Python sobre Raspberry Pi.

\section*{B.1 — Interfaces abstractas del núcleo}

\begin{lstlisting}[language=Python, caption={Interfaz ITransport — contrato de la capa de transporte}]
from __future__ import annotations
from abc import ABC, abstractmethod
from types import TracebackType


class ITransport(ABC):

    @abstractmethod
    def connect(self) -> None: ...

    @abstractmethod
    def disconnect(self) -> None: ...

    @abstractmethod
    def send(self, payload: bytes) -> None: ...

    @abstractmethod
    def receive(self) -> bytes: ...

    @abstractmethod
    def __enter__(self) -> ITransport: ...

    @abstractmethod
    def __exit__(
        self,
        exc_type: type[BaseException] | None,
        exc_val: BaseException | None,
        exc_tb: TracebackType | None,
    ) -> bool | None: ...
\end{lstlisting}

\section*{B.2 — Transporte ISO-TP sobre SocketCAN}

\begin{lstlisting}[language=Python, caption={IsoTpTransport — implementacion concreta sobre python-can + can-isotp}]
import time
import can
import isotp
from config.can_config import CAN_RX_ID, CAN_TX_ID, ISOTP_CF_SEPARATION_MS, ISOTP_PADDING_BYTE
from core.exceptions import DiagnosticTimeoutError
from core.interfaces.i_transport import ITransport


class IsoTpTransport(ITransport):

    def __init__(self, channel: str = "can0", timeout: float = 5.0) -> None:
        self._channel = channel
        self._timeout = timeout
        self._address = isotp.Address(
            addressing_mode=isotp.AddressingMode.Normal_11bits,
            txid=CAN_TX_ID,
            rxid=CAN_RX_ID,
        )
        self._params = {
            "stmin": ISOTP_CF_SEPARATION_MS,
            "blocksize": 0,
            "tx_padding": ISOTP_PADDING_BYTE,
        }
        self._stack: isotp.CanStack | None = None
        self._bus: can.BusABC | None = None

    def connect(self) -> None:
        self._bus = can.Bus(channel=self._channel, interface="socketcan")
        # Vaciar tramas residuales del buffer del kernel
        while self._bus.recv(timeout=0) is not None:
            pass
        self._stack = isotp.CanStack(self._bus, address=self._address,
                                     params=self._params)
        self._flush_rx()

    def disconnect(self) -> None:
        if self._stack is None:
            return
        self._bus.shutdown()
        self._stack = None
        self._bus = None

    def send(self, payload: bytes) -> None:
        self._flush_rx()
        self._stack.send(payload)
        while self._stack.transmitting():
            self._stack.process()
            time.sleep(0.001)

    def receive(self) -> bytes:
        deadline = time.monotonic() + self._timeout
        while time.monotonic() < deadline:
            self._stack.process()
            if self._stack.available():
                return self._stack.recv()
            time.sleep(0.001)
        raise DiagnosticTimeoutError(f"No response within {self._timeout:.1f}s.")

    def _flush_rx(self) -> None:
        self._stack.process()
        while self._stack.available():
            self._stack.recv()
            self._stack.process()
\end{lstlisting}

\section*{B.3 — Logger SQLite con modo WAL}

\begin{lstlisting}[language=Python, caption={SqliteDataLogger — esquema, buffer de escritura y modo WAL}]
import sqlite3, threading
from datetime import datetime, timezone
from core.interfaces.i_data_logger import IDataLogger

_SCHEMA = """
PRAGMA journal_mode=WAL;

CREATE TABLE IF NOT EXISTS sessions (
    id         INTEGER PRIMARY KEY AUTOINCREMENT,
    label      TEXT    NOT NULL DEFAULT '',
    started_at TEXT    NOT NULL,
    ended_at   TEXT
);

CREATE TABLE IF NOT EXISTS samples (
    id           INTEGER PRIMARY KEY AUTOINCREMENT,
    session_id   INTEGER NOT NULL REFERENCES sessions(id),
    pid          INTEGER NOT NULL,
    name         TEXT    NOT NULL,
    value        REAL    NOT NULL,
    unit         TEXT    NOT NULL,
    monotonic_ts REAL    NOT NULL,
    wall_ts      TEXT    NOT NULL
);

CREATE TABLE IF NOT EXISTS commands (
    id           INTEGER PRIMARY KEY AUTOINCREMENT,
    session_id   INTEGER NOT NULL REFERENCES sessions(id),
    command      TEXT    NOT NULL,
    request_hex  TEXT    NOT NULL,
    response_hex TEXT    NOT NULL,
    wall_ts      TEXT    NOT NULL
);
"""

_BUFFER_SIZE = 50


class SqliteDataLogger(IDataLogger):
    # WAL permite escrituras concurrentes desde hilo monitor y hilo BLE.

    def __init__(self, db_path: str = "diagnostics.db") -> None:
        self._lock = threading.Lock()
        self._sample_buffer: list[tuple] = []
        self._conn = sqlite3.connect(db_path, check_same_thread=False)
        self._conn.executescript(_SCHEMA)
        self._conn.commit()

    def log_sample(self, session_id: int, sample) -> None:
        row = (session_id, sample.pid, sample.name, sample.value,
               sample.unit, sample.timestamp,
               datetime.now(timezone.utc).isoformat())
        with self._lock:
            self._sample_buffer.append(row)
            if len(self._sample_buffer) >= _BUFFER_SIZE:
                self._flush_buffer()

    def log_command(self, session_id, command, request, response) -> None:
        with self._lock:
            self._conn.execute(
                "INSERT INTO commands"
                " (session_id, command, request_hex, response_hex, wall_ts)"
                " VALUES (?, ?, ?, ?, ?)",
                (session_id, command, request.hex(), response.hex(),
                 datetime.now(timezone.utc).isoformat()),
            )
            self._conn.commit()

    def _flush_buffer(self) -> None:
        if not self._sample_buffer:
            return
        self._conn.executemany(
            "INSERT INTO samples"
            " (session_id, pid, name, value, unit, monotonic_ts, wall_ts)"
            " VALUES (?, ?, ?, ?, ?, ?, ?)",
            self._sample_buffer,
        )
        self._conn.commit()
        self._sample_buffer.clear()
\end{lstlisting}

\section*{B.4 — Decorador LoggingTransport}

\begin{lstlisting}[language=Python, caption={LoggingTransport — patron decorador sobre ITransport}]
from core.interfaces.i_transport import ITransport
from core.interfaces.i_data_logger import IDataLogger


class LoggingTransport(ITransport):
    """Envuelve cualquier ITransport e intercepta send/receive para loggear
    las tramas CAN en hexadecimal en SQLite sin modificar el transporte base."""

    def __init__(self, inner: ITransport, logger: IDataLogger,
                 session_id: int, command: str) -> None:
        self._inner = inner
        self._logger = logger
        self._session_id = session_id
        self._command = command
        self._last_sent: bytes = b""

    def connect(self) -> None:    self._inner.connect()
    def disconnect(self) -> None: self._inner.disconnect()
    def __enter__(self):          self._inner.connect(); return self
    def __exit__(self, *a):       self._inner.disconnect()

    def send(self, payload: bytes) -> None:
        self._last_sent = payload
        self._inner.send(payload)

    def receive(self) -> bytes:
        response = self._inner.receive()
        self._logger.log_command(
            self._session_id, self._command,
            self._last_sent, response,
        )
        return response
\end{lstlisting}
